\documentclass[11pt,a4paper]{article}
\usepackage{fontspec}
\usepackage[turkish,shorthands=off]{babel}
\usepackage[margin=2.2cm]{geometry}
\usepackage{booktabs}
\usepackage{longtable}
\usepackage{array}
\usepackage{amsmath}
\usepackage{microtype}
\usepackage{enumitem}
\usepackage[colorlinks=true,linkcolor=blue,urlcolor=blue,citecolor=blue]{hyperref}

\setlength{\parskip}{4pt}
\setlength{\parindent}{0pt}

\newcommand{\etal}{\emph{vd.}}

\title{İnsansız Hava Aracı Telemetrisi ve ADS-B Hava Trafiği Verisinde\\
Anomali Tespiti: Yöntemler, Deneysel Sonuçlar ve\\
Kapsamlı Başarısızlık Analizi}
\author{Anomali Tespiti Projesi --- Neden--Sonuç Raporu}
\date{17 Temmuz 2026}

\begin{document}
\maketitle

\begin{abstract}
Bu rapor, insansız hava aracı (İHA) uçuş telemetrisi ve geniş alanlı ADS-B hava
trafiği verisi üzerinde yürütülen anomali tespiti projesinin neden hedeflenen
operasyonel başarıya ulaşamadığını, veri seti, kolon, yöntem ve ölçülmüş skor
düzeyinde belgeler. Beş etiketli İHA veri seti ve 2.76 milyar satırlık etiketsiz
bir ADS-B külliyatı üzerinde dokuz yöntem ailesi (Isolation Forest, gradyan
artırma, yoğun ve LSTM otokodlayıcılar, USAD, LSTM öngörücü, sıfır-atış temel
model, kural-bazlı fizik artıkları ve bağlamsal olasılıksal LSTM; karar
katmanında CUSUM ve conformal kalibrasyon) denenmiştir. Araştırma metriklerinde
literatürle uyumlu sonuçlar elde edilmiştir (örn.\ ALFA veri setinde LSTM-AE ile
AUPRC 0.872); buna karşın \textbf{hiçbir konfigürasyon, dondurulmuş eşikle
birlikte ölçülen olay-yakalama + saatlik yanlış alarm bütçesi biçimindeki
operasyonel kabul ölçütünü hiçbir veri setinde karşılayamamıştır}. Başarısızlığın
kök nedenleri; değerlendirme çıtası uyumsuzluğu, yapısal olarak yetersiz normal
uçuş maruziyeti, heterojen normal havuzlarında yenilik-tespitinin çökmesi,
yeniden-kurma tabanlı modellerdeki genlik-baskınlığı artefaktı, etiket
granülerliği uyumsuzluğu, fiziksel gözlenebilirlik boşlukları ve alan (domain)
kayması olarak tespit edilmiş ve her biri ölçülmüş sayılarla kanıtlanmıştır.
Kritik ayrım şudur: son iki çalışmada sinyalin \emph{varlığı} eşikten bağımsız
istatistiksel testlerle kanıtlanmış, fakat operasyonel eşik, kalibrasyon için
gereken bağımsız normal uçuş-saatinin veri setlerinde mevcut olmaması nedeniyle
kurulamamıştır. Yani sonuç ``sinyal yok'' değil, ``sinyal var; mevcut veriyle
operasyonel karar sınırı kalibre edilemez'' biçimindedir.
\end{abstract}

\tableofcontents
\newpage

\section{Yönetici Özeti}

Projenin amacı, İHA uçuş telemetrisinde (sensör/aktüatör arızaları, GNSS
bütünlük bozulmaları, siber saldırılar) ve ticari hava trafiği yayınlarında
(ADS-B) anomalileri, operasyonel olarak kabul edilebilir bir yanlış alarm
yüküyle tespit eden bir sistem kurmaktı. Sonuç, tek cümleyle:

\begin{quote}
\textbf{Denenen hiçbir yöntem--veri seti kombinasyonu, ``dondurulmuş tek bir
eşikle, aynı anda ölçülen olay-yakalama oranı ve saatlik yanlış alarm bütçesi''
biçimindeki operasyonel kabul ölçütünü karşılayamadı.} Buna karşılık,
literatürün raporladığı eşikten-bağımsız sıralama metriklerinde (AUROC/AUPRC)
sonuçlar yayınlanmış çalışmalarla aynı ligdedir.
\end{quote}

Bu iki cümlenin birlikte okunması raporun ana tezidir: başarısızlık, yöntemlerin
literatürdeki hâllerinin ``çalıştırılamaması'' değildir. Literatürün genellikle
\emph{ölçmediği} bir büyüklüğe --- dondurulmuş eşikte saatlik yanlış alarm
yüküne --- hedef konmuş; küçük ve heterojen akademik veri setleri bu hedefi
istatistiksel olarak taşıyamamıştır. Nihai iki çalışma (GNSS bütünlük pilotu ve
komut$\to$tepki artık modeli) disiplinli birer NO-GO kararıyla kapanmış; her
ikisinde de kör test kümeleri hiç açılmamış, eşikler sonuca bakılarak
gevşetilmemiştir.

Sayısal özet:

\begin{itemize}[itemsep=2pt]
  \item 5 etiketli İHA veri seti (ALFA, UAV Attack, UAV-SEAD, RflyMAD ve
        RflyMAD'in GNSS alt-kümesi) + 1 etiketsiz ADS-B külliyatı
        (2.762.067.817 satır, 386.779 uçak, $\sim$11 ay).
  \item 9 yöntem ailesi, $\sim$50'nin üzerinde konfigürasyon.
  \item Araştırma metriği örneği: ALFA'da LSTM-AE ile AUPRC \textbf{0.872}
        (Isolation Forest ile 0.858, LightGBM ile 0.843); literatürle uyumlu
        düzey.
  \item Operasyonel metrik örneği: en iyi tekil kanal (itki komutu artığı,
        RflyMAD) \textbf{0.526} olay-yakalama sağladı ama \textbf{22.28 yanlış
        alarm/uçuş-saati} üretti; bütçe 12 idi. Kritik bütçe (2/saat) hiçbir
        yerde yaklaşılamadı.
  \item Nihai kalibrasyon analizi: yanlış alarm bütçesini güvenilir çözmek için
        gereken bağımsız normal uçuş-saatinin yalnızca $1/11.8$'i (sabit kanat)
        ve $1/5.1$'i (çok-rotor) veri setlerinde mevcuttur; resmî külliyat
        tavanları bu açığı matematiksel olarak kapatamamaktadır.
\end{itemize}

\section{Problem Tanımı ve Başarı Ölçütleri}
\label{sec:olcut}

\subsection{İki metrik ailesi}

Zaman serisi anomali tespitinde iki farklı soru vardır ve proje boyunca ikisi
bilinçli olarak ayrı tutulmuştur:

\begin{enumerate}[itemsep=2pt]
  \item \textbf{Araştırma sorusu (eşikten bağımsız sıralama):} skor, anomalili
        örnekleri normallerden ne kadar iyi ayırıyor? Metrikler: AUROC, AUPRC.
        Literatürdeki ``\%95 başarı'' raporlarının büyük kısmı bu ailedendir.
  \item \textbf{Operasyonel soru (dondurulmuş eşikte eşzamanlı çift metrik):}
        eşik, sonuç görülmeden \emph{önce} yalnız normal veriden kalibre edilip
        dondurulduğunda, (a) etiketli arıza olaylarının ne kadarı zamanında
        yakalanıyor ve (b) normal uçuşlarda saatte kaç yanlış alarm üretiliyor?
        Bu iki sayı \emph{aynı eşikle, aynı anda} raporlanır; asla ayrı ayrı
        değil.
\end{enumerate}

Projenin kabul ölçütleri baştan şöyle donduruldu (sonuç görülmeden):

\begin{center}
\small
\begin{tabular}{lccc}
\toprule
Sözleşme & Olay-yakalama (recall) & Yanlış alarm bütçesi & Gecikme hedefi \\
\midrule
Kritik   & $\geq 0.30$ & $\leq 2$ alarm/uçuş-saati  & $\leq 5$ s \\
Uyarı (advisory) & $\geq 0.50$ & $\leq 12$ alarm/uçuş-saati & $\leq 15$ s \\
\bottomrule
\end{tabular}
\end{center}

\subsection{Neden bu ölçüt? Literatürdeki ölçütün bilinen kusurları}

Bu seçim keyfî değildir; zaman serisi anomali tespiti değerlendirme
literatürünün son yıllardaki iki temel eleştirisine dayanır:

\begin{itemize}[itemsep=2pt]
  \item \textbf{Point-adjust şişirmesi:} olay içinde tek bir noktayı yakalamayı
        ``bütün olay yakalandı'' saymak, rastgele skorları bile başarılı
        gösterebilir \cite{kim2022rigorous}. Projede bu düzeltme kendi
        sonuçlarımıza uygulandığında olay-başlangıcı yakalama oranı
        $0.594 \to 0.194\text{--}0.224$'e düşmüştür (\S\ref{sec:alfa-sonuc}).
  \item \textbf{Kusurlu kıyas kümeleri:} yaygın akademik kıyas kümelerinin
        önemsiz (trivial) çözülebilirlik ve etiket hataları içerdiği
        gösterilmiştir \cite{wu2021flawed}. Kapsamlı karşılaştırmalar, tek bir
        yöntemin genel üstünlüğü olmadığını bulmaktadır \cite{schmidl2022eval}.
\end{itemize}

Saatlik yanlış alarm bütçesi ise operasyonel izleme pratiğinin standart
büyüklüğüdür: alarm üreten ama operatörünü alarm yağmuruna boğan bir dedektör
sahada kapatılır. Raporun geri kalanında görüleceği gibi, \emph{proje dört
etiketli veri setinde de hep aynı duvara --- bu bütçeye --- toslamıştır};
eşikten bağımsız metrikler çoğu zaman makul, hatta iyiydi.

\section{Veri Setleri: İçerik, Kolon Envanteri ve Kolon-Bazında Uygulanan
Operasyonlar}
\label{sec:veri}

Tablo~\ref{tab:veri-ozet} kullanılan altı veri kaynağını özetler. Kolon
bileşimleri ağırlıkla \textbf{sensör verisi}dir: yüksek frekanslı ataletsel
ölçümler (jiroskop, ivmeölçer), otopilot komut/ölçüm çiftleri, aktüatör PWM
çıkışları, GNSS konum/hız ve --- ADS-B tarafında --- uçağın kendi yayınladığı
kinematik durum vektörü ile bütünlük meta-verisi.

\begin{table}[h]
\centering\small
\caption{Veri setlerinin özeti.}
\label{tab:veri-ozet}
\begin{tabular}{p{2.4cm}p{2.6cm}p{3.4cm}p{2.9cm}p{3.2cm}}
\toprule
Veri seti & Platform & Anomali türleri & Hacim & Normal havuz \\
\midrule
ALFA \cite{keipour2021alfa} & Sabit kanat İHA, gerçek uçuş & Motor durması,
kumanda yüzeyi (rudder, elevator, aileron) arızaları & 47 uçuş (sabit resmî
külliyat) & \textbf{11 normal uçuş} \\
\addlinespace[2pt]
UAV Attack \cite{whelan2020uav} & PX4 çok-rotor & GPS spoofing, GPS jamming,
ağ-katmanı DoS & $\sim$78 bin satır & \textbf{6 normal log} \\
\addlinespace[2pt]
UAV-SEAD & Çok-rotor, ağırlıkla simülasyon & GPS spoofing; irtifa, mekanik,
konum anomalileri & $\sim$1396 uçuş & 398 normal $\approx$ \textbf{64 bağımsız
oturum} (sonradan 899 normale büyütüldü) \\
\addlinespace[2pt]
RflyMAD \cite{rflymad} & Çok-rotor; gerçek + SIL/HIL & Motor ve sensör
arızaları & 490 gerçek uçuş (242 motor, 197 sensör) & \textbf{51 arızasız uçuş}
(resmî tavan 84) \\
\addlinespace[2pt]
ADS-B (adsb.lol) \cite{adsblol} & Ticari hava trafiği & Etiketsiz; sentetik
bozulma senaryolarıyla değerlendirme & \textbf{2.76 milyar satır}, 386.779
uçak, $\sim$11 ay & Etiketsiz (doğal trafik) \\
\addlinespace[2pt]
GNSS alt-çalışması & RflyMAD GPS alt-kümesi + SIL/HIL rüzgâr & GNSS gürültü ve
ölçek-faktörü arızaları & 94 uçuş (rol ayrımlı) + 886 rüzgâr vakası & Fit 20 +
kalibrasyon 10 uçuş \\
\bottomrule
\end{tabular}
\end{table}

Aşağıdaki alt bölümler, her veri setinin kolonlarını ve \emph{her kolon grubuna
uygulanan somut operasyonları} verir. Bu envanter raporun en önemli
bölümlerinden biridir, çünkü başarısızlık nedenlerinin birkaçı doğrudan kolon
düzeyindeki özelliklerden (örneklem hızı uyuşmazlığı, sıfır-varyanslı alanlar,
etiket granülerliği) doğmuştur.

\subsection{ALFA --- sabit kanat İHA arıza veri seti}
\label{sec:alfa-veri}

ALFA \cite{keipour2021alfa}, Carnegie Mellon Üniversitesi'nin yayımladığı,
gerçek uçuşta kontrollü arıza enjeksiyonu içeren 47 uçuşluk sabit bir
külliyattır. Sınıf dağılımı ciddi biçimde dengesizdir: motor arızası çoğunluk
sınıfıdır; rudder $n{=}4$, elevator $n{=}2$, aileron+rudder $n{=}1$ ve yalnız
\textbf{11 normal uçuş} vardır. Kolonlar tamamen uçuş kontrol sensörü
ağırlıklıdır:

\begin{table}[h]
\centering\small
\caption{ALFA kolon envanteri ve uygulanan operasyonlar.}
\label{tab:alfa-kolon}
\begin{tabular}{p{3.6cm}p{3.1cm}p{1cm}p{6.6cm}}
\toprule
Kolon grubu & Fiziksel anlam & $\sim$Hz & Uygulanan operasyonlar \\
\midrule
Otopilot komut/ölçüm çiftleri (istenen ve ölçülen roll, pitch, yaw, hava hızı) &
Komut$\to$tepki artıklarının ana hammaddesi & 20--25 & Açı farkları sarım-duyarlı
($\operatorname{atan2}$) hesaplandı; komut--tepki artığı türetildi; medyan/MAD
ile sağlam ölçekleme + kırpma \\
\addlinespace[2pt]
Ataletsel ölçümler (jiroskop, ivmeölçer, kuaterniyon) & En yüksek hızlı gerçek
dinamik; referans saat & 45--50 & Kuaterniyon işaret sürekliliği düzeltmesi
($q \leftrightarrow -q$ aynı duruş); ivme türevlenmedi, bağlam olarak kullanıldı \\
\addlinespace[2pt]
Aktüatör PWM çıkışları (servo, gaz) & Fiilî aktüatör komutu & 20 & Uçuşlar arası
kalibrasyon kayması nedeniyle uçuş-içi trim medyanına göre normalize edildi \\
\addlinespace[2pt]
Küresel/yerel konum, yer hızı, irtifa & Konum/kinematik & 5--10 & GPS irtifa
sıçramalarına karşı aralık kapıları; yer hızı ile hava hızı ayrı kanallar olarak
tutuldu (rüzgâr nedeniyle asla birbirinin yerine kullanılmadı) \\
\addlinespace[2pt]
Yörünge sapması (çapraz-iz hatası) & Planlanan rotadan yanal sapma & 10 &
Nedensel ölçümde en güçlü tekil kanal; yol-noktası geçişindeki yapısal sıçrama
için $\pm 2$ s maske; V-dönüşü olayları için altı-parametreli maske (32 uçuşta 9
olay, 697 satır maskelendi) \\
\addlinespace[2pt]
Arıza durumu (yalnız arıza aktifken yayınlanır) & Etiket & olay & Başlangıç
(onset) = ilk yayın; öncesindeki 10 s koruma bandı eğitimden dışlandı \\
\bottomrule
\end{tabular}
\end{table}

\textbf{Hijyen kuralları (tüm kolonlara):} $\Delta t \leq 0$ satırlar atıldı ve
sayıldı; $\Delta t > 5\times$medyan boşluk olarak işaretlendi; 2 s boyunca
değişmeyen yüksek-hızlı kanal ``bayat'' (stale) bayrağı aldı ve o pencerede
artık hesaplanmadı; fiziksel aralık ihlali ($|\text{roll}| \leq 180^\circ$,
hava hızı 0--60 m/s vb.) olan uçuşlar karantinaya alındı; sessiz kırpma
yapılmadı.

\textbf{Kritik metodolojik dönüş (interpolasyon):} projenin ilk turlarında tüm
kanallar 50 Hz'e doğrusal interpole edilmişti. Bu, iki zarar üretti: (i) 5 Hz'lik
GNSS kanalını 50 Hz'e interpole etmek ardışık 10 örneği yapay bir doğru üzerine
dizer ve model bu yapay pürüzsüzlüğü öğrenir; (ii) veri kesintisi sırasında
interpolant iki uzak nokta arasında köprü kurarak tam anomali anını
pürüzsüzleştirir. Son turda interpolasyon tamamen kaldırıldı; her kanal kendi
doğal hızında tutulup hizalama geçmişe-bakan asof birleştirmeyle
(\texttt{merge\_asof}, backward) yapıldı ve her düşük hızlı kanala bir bayatlık
(staleness) kolonu eşlik etti.

\subsection{UAV Attack --- çok-rotor siber saldırı veri seti}

Whelan \etal{} \cite{whelan2020uav} tarafından yayımlanan bu set, PX4 tabanlı
çok-rotorlara GPS spoofing, GPS jamming ve ağ-katmanı DoS saldırıları içerir.
Kolonlar PX4 telemetrisidir (konum, hız, ataletsel ölçümler, EKF durumları).
Uygulanan operasyonlar: kanal bazında fiziksel artık türetimi (örn.\ GPS hız
artığı), sağlam ölçekleme, pencere özetleri.

Bu setin iki yapısal sınırı ölçümle kanıtlandı:

\begin{itemize}[itemsep=2pt]
  \item \textbf{Gözlenebilirlik boşluğu:} ağ-katmanı DoS saldırısı fiziksel
        telemetriye \emph{hiç yansımıyor} --- 6 saldırı logundan 4'ü hiçbir
        kanalda tespit edilemedi. Ağ katmanındaki bir imzayı uçuş fiziği
        kolonlarında aramak kategori hatasıdır; bu, yöntem değil veri--problem
        eşleşmesi başarısızlığıdır.
  \item \textbf{Sinsi saldırı ölçeği:} 2 m/s'lik kademeli GPS kayması, GPS hız
        artığını ancak normal doğrulama maksimumu düzeyinde oynatmaktadır;
        yani bu şiddette spoofing, artık uzayında normal varyasyondan
        ayrılamamaktadır.
  \item \textbf{Normal havuzu 6 log:} saatlik yanlış alarm bütçesi gibi bir
        büyüklüğü 6 normal kayıtla kalibre etmek istatistiksel olarak
        anlamsızdır.
\end{itemize}

\subsection{UAV-SEAD --- simülasyon ağırlıklı çok-rotor anomali seti}

$\sim$1396 uçuşluk, GPS spoofing ile irtifa/mekanik/konum anomalileri içeren,
ağırlıkla simülasyon kaynaklı bir settir. Görünürde 398 normal uçuş vardır;
ancak oturum analizi bunların yalnız \textbf{$\sim$64 bağımsız oturumdan}
geldiğini göstermiştir --- aynı oturumun uçuşları güçlü biçimde koreledir ve
etkin örneklem, görünenden çok daha küçüktür. Kolonlar: konum (X, Y, Z),
aktüatör komutları (motor itki komutları dahil), duruş ve türetilmiş kinematik.

Uygulanan operasyonlar: pencere-bazlı özet özellikler (ilk turlarda 73--85
özellik); kategoriye özgü fiziksel kanallar (dikey tutarlılık: konum-Z ile
tırmanma dinamiği; motor simetrisi: aktüatör çıkışları arası denge; itki
komutu); CUSUM ve K-of-N kalıcılık karar katmanları; oturum-bazlı bölme ve 5
tohumlu tekrar. Sonradan normal havuz 899 uçuşa büyütüldü; \S\ref{sec:sead-sonuc}'da
görüleceği gibi bu büyütme operasyonel sonucu kurtarmadı.

\subsection{RflyMAD --- gerçek çok-rotor arıza veri seti}
\label{sec:rfly-veri}

RflyMAD \cite{rflymad}, Beihang Üniversitesi'nin çok-rotor arıza veri setidir;
projede yalnız \emph{gerçek} uçuş alt-kümesi kullanıldı: 490 uçuş (242 motor
arızası, 197 sensör arızası, 51 arızasız; resmî sayfada arızasız tavanı 84).
Kayıtlar PX4 ULog biçimindedir:

\begin{table}[h]
\centering\small
\caption{RflyMAD kolon envanteri ve uygulanan operasyonlar.}
\label{tab:rfly-kolon}
\begin{tabular}{p{3.6cm}p{3.6cm}p{7cm}}
\toprule
Kolon grubu (ULog konusu) & Fiziksel anlam & Uygulanan operasyonlar \\
\midrule
Aktüatör çıkışları (4 motor PWM) & Fiilî motor komutları & Motorlar arası
dağılım/asimetri kanalı türetildi (toplam itki ve duruş komutuna koşullu);
medyan/MAD ölçekleme \\
\addlinespace[2pt]
Duruş + duruş hedefi (setpoint) & Komut$\to$tepki çifti & Hedef--ölçüm artığı;
sarım-duyarlı açı farkları \\
\addlinespace[2pt]
Yerel konum/hız & Kinematik & Konum hedefi$\to$hız tepkisi kanalı; geçmişe-bakan
asof hizalama, ileri-doldurma YOK \\
\addlinespace[2pt]
Batarya durumu & Voltaj çökmesi motor arızasını taklit edebilir & Yalnız bağlam
değişkeni olarak alındı; dedektör kanalı yapılmadı \\
\addlinespace[2pt]
ESC durumu (RPM, akım, voltaj, sıcaklık) & Motor sağlığının doğrudan ölçümü &
\textbf{Kullanılamadı:} incelenen tüm uçuşlarda bu alanlar sabit sıfır ---
şema var, veri yok. Motor-sağlığı pilotu bu nedenle hiç açılmadı (gözlenebilirlik
sözleşmesi ihlal edilmeden yürütülemezdi) \\
\addlinespace[2pt]
Arıza aralığı gerçek-referansı & Olay bazlı etiket (başlangıç/bitiş) & Uçuş-düzeyi
vekil etiket yerine aralık etiketi kullanıldı; klasör etiketiyle çelişen 5 uçuş
ve yanlış arıza kimliği taşıyan 6 kayıt karantinaya alındı \\
\bottomrule
\end{tabular}
\end{table}

\textbf{Bu setteki en öğretici düzeltme} etiket granülerliğidir: ilk değerlendirme
``uçuşun tamamı arızalı'' vekil etiketiyle 0.749 olay-yakalama gösteriyordu.
Gerçek arıza \emph{aralığı} etiketine geçilince aynı yöntem
$0.526$'ya düştü (\S\ref{sec:rfly-sonuc}). Kolay soruya verilmiş doğru cevap,
zor sorunun cevabı değildir.

\subsection{ADS-B (adsb.lol) --- geniş alanlı hava trafiği külliyatı}
\label{sec:adsb-veri}

adsb.lol \cite{adsblol} topluluk ağının günlük tam-ağ arşivlerinden 30 anlık
görüntü ($\sim$11 aylık pencere) işlendi: \textbf{2.762.067.817 satır, 386.779
benzersiz uçak}. Veri, uçakların kendi yayınladığı durum vektörüdür; kolon
bileşimi kinematik-ağırlıklıdır ve bütünlük meta-verisi içerir:

\begin{table}[h]
\centering\small
\caption{ADS-B kolon envanteri ve uygulanan operasyonlar.}
\label{tab:adsb-kolon}
\begin{tabular}{p{3.4cm}p{3.4cm}p{7.4cm}}
\toprule
Kolon & Anlam & Uygulanan operasyonlar \\
\midrule
Zaman damgası & Mesaj zamanı & $\Delta t$ hesaplandı; $\log(1+\Delta t)$ model
girdisi; ardışık mesaj arası $> 60$ s ise pencere üretimi kesildi; $\Delta t
\leq 0$ atıldı \\
\addlinespace[2pt]
Enlem/boylam & Konum & Ardışık konumdan hız vektörü ve büyük-daire kerterizi
türetildi; hız/yön artıklarının referansı \\
\addlinespace[2pt]
Barometrik irtifa & İrtifa & $\Delta\text{irtifa}/\Delta t$ türetildi (dikey hız
artığının referansı); \texttt{"ground"} değeri sayı değil bayrağa çevrildi \\
\addlinespace[2pt]
Geometrik irtifa & GNSS irtifası & Barometrik--geometrik fark türevi kanal adayıydı;
fit verisinde MAD tam sıfır çıktığı için \emph{yapay taban verilmeden dışlandı}
(aktif kanal sayısı 6 değil 5 kaldı) \\
\addlinespace[2pt]
Yer hızı & Yayınlanan hız & Hız artığı = yayınlanan hız $-$ konumdan türetilen
hız \\
\addlinespace[2pt]
Yer izi (track) & Yayınlanan yön & $\sin/\cos$ kodlaması (dairesel süreklilik);
yön artığı = yayınlanan track $-$ konum kerterizi (sarım-duyarlı) \\
\addlinespace[2pt]
Dikey hız & Yayınlanan tırmanma/alçalma & Dikey hız artığı = yayınlanan değer $-$
$\Delta\text{irtifa}/\Delta t$ \\
\addlinespace[2pt]
Doğu/kuzey hız bileşenleri & Hız vektörü & Bileşen bazında hız-vektörü artıkları
(sinsi konum kayması için birincil kanıt kanalı) \\
\addlinespace[2pt]
Squawk (transponder kodu) & Pilot beyanı & Kural kanalı: acil kodlar
(7500/7600/7700) doğrudan, öğrenmesiz sinyal \\
\addlinespace[2pt]
NIC / NACp / SIL & Yayınlanan bütünlük/doğruluk kategorileri & Kural-bazlı
``beyan edilen kalite düşüşü'' sinyali; sinir ağına karıştırılmadı \\
\addlinespace[2pt]
Uçuş fazı, kadans & Türetilmiş bağlam & Faz yalnız geçmiş satırlardan çıkarıldı
(sızıntısız); kadans kovaları girdiye eklendi; her girdiye kullanılabilirlik
maskesi eşlik etti \\
\bottomrule
\end{tabular}
\end{table}

Tüm artık kanalları, fit-normal medyan/MAD ile ölçeklenip $[-5,5]$'te kırpıldı:
\[
z = \operatorname{clip}\!\big((r - \operatorname{med})/\operatorname{MAD},\,
-5,\, 5\big).
\]
Pencereler uçuş sınırını asla geçmedi.

\textbf{Etiketsizlik ve değerlendirme:} gerçek anomali etiketi olmadığından
değerlendirme, doğal veriye enjekte edilen beş sentetik bozulma senaryosuyla
yapıldı: yer hızı sapması (bias), dikey hız donması, track donması, sinsi konum
rampası, irtifa kesintisi. Ayrıca pencere-düzeyi vekil etiketlemenin yapısal bir
AUC tavanı ($\sim$0.75) olduğu ölçüldü: mükemmel bir dedektör bile bu
etiketlemeyle 0.75 civarında görünür. Bu tavan, sonuçların yorumunda hep dikkate
alındı.

\subsection{GNSS bütünlük alt-çalışması (RflyMAD GPS alt-kümesi)}
\label{sec:gnss-veri}

Projenin sondan bir önceki çalışması, dar ve ön-kayıtlı bir soruya indirgendi:
``RflyMAD gerçek uçuşlarındaki GNSS gürültü ve ölçek-faktörü arızaları, kabul
edilebilir operatör yüküyle zamanında tespit edilebilir mi?'' Veri rolleri
baştan ayrıldı ve hiç ihlal edilmedi: fit 20, kalibrasyon 10, geliştirme 23,
prova 15, \emph{mühürlü kör test 20} (hiç açılmadı), karantina 6 uçuş.

Kolonlar ve operasyonlar:

\begin{itemize}[itemsep=2pt]
  \item \textbf{EKF yenilik (innovation) değerleri} ($\sim$2 Hz temel eksen):
        GPS yatay konum, yatay hız ve dikey hız yenilikleri, ilgili yenilik
        varyansının kareköküne bölünerek normalize edildi --- literatürdeki
        model-bazlı arıza tespitinin (FDI) standart büyüklüğü
        \cite{willsky1976survey,isermann2006fdi}.
  \item \textbf{EKF test oranı, bayrak/sıfırlama durumları; GPS fix tipi, yatay/dikey
        konum belirsizliği (EPh/EPv), HDOP/VDOP, uydu sayısı, gürültü ve jamming
        göstergeleri; yerel hız:} geçmişe-bakan asof birleştirmeyle eklendi;
        \emph{ileri doldurma yapılmadı} (arıza anını normal veriyle örtmemek
        için).
  \item Alan-kayması stres testi için 886 SIL/HIL rüzgâr vakası ayrıca işlendi;
        bunlar GNSS arızası içermez, yalnız dondurulmuş dedektörlerin çevresel
        bozucudaki alarm yükünü ölçer.
\end{itemize}

\section{Yöntemler}
\label{sec:yontem}

Denenen yöntemler, literatürdeki dört ana aileyi kapsar. Hepsi ``yalnız
normalden öğren, sapanı işaretle'' (novelty detection) veya zayıf-gözetimli
kurulumla çalıştırıldı; kör test kümelerine hiçbir aşamada bakılmadı.

\subsection{Sıra dışılık skorlayıcıları}

\begin{enumerate}[itemsep=3pt]
  \item \textbf{Isolation Forest} \cite{liu2008iforest}: satır- ve
        pencere-düzeyi izolasyon derinliği skoru; hem monolitik hem kanal-bazlı
        (modüler füzyon) varyantları.
  \item \textbf{Gradyan artırma (LightGBM)} \cite{ke2017lightgbm}:
        pencere özellikleri üzerinde zayıf-gözetimli sınıflandırma.
  \item \textbf{Yoğun (dense) otokodlayıcı}: yeniden-kurma hatası skoru.
  \item \textbf{LSTM otokodlayıcı} \cite{malhotra2016lstm}: sıralı
        yeniden-kurma; literatürde çok-sensörlü telemetri için standart temel
        yöntem.
  \item \textbf{USAD} \cite{audibert2020usad}: çekişmeli (adversarial) eğitimli
        otokodlayıcı.
  \item \textbf{LSTM öngörücü (forecaster)}: bir-adım-ileri tahmin hatası
        skoru.
  \item \textbf{Sıfır-atış temel model (Chronos)} \cite{ansari2024chronos}:
        önceden eğitilmiş zaman serisi temel modelinin tahmin artıkları;
        alan-içi eğitim yapılmadan kullanıldı.
  \item \textbf{Kural-bazlı fizik tutarlılık skorlayıcı}: \S\ref{sec:adsb-veri}'deki
        fiziksel artıklar üzerinde el yapımı ceza kuralları; öğrenme yok.
  \item \textbf{Bağlamsal olasılıksal LSTM}: ham sinyali yeniden kurmak yerine,
        uçuşun önceki 12 satırından her artık kanalı için beklenen merkez
        $\mu$ ve normal oynaklık $\sigma$ tahmin eder; eğitim kaybı maskeli
        Gauss negatif log-olabilirliğidir:
        \[
        \mathrm{NLL}_{c,t} = \log \sigma_{c,t} +
        \tfrac{1}{2}\Big(\tfrac{z_{c,t}-\mu_{c,t}}{\sigma_{c,t}}\Big)^{2},
        \qquad
        S_{c,t} = \frac{|z_{c,t}-\mu_{c,t}|}{\sigma_{c,t}} .
        \]
        Skor $S$, kanal/faz/kadans koşullu conformal $p$-değerine
        \cite{vovk2005conformal} çevrilecek şekilde tasarlandı. Tepki
        değişkeninin kendi kısa geçmişi girdiye \emph{bilinçli olarak}
        verilmedi (otoregresif kopyacılık tuzağına karşı, \S\ref{sec:kn-artefakt}).
  \item \textbf{Komut$\to$tepki artık modeli (model-bazlı FDI çerçevesi)}
        \cite{willsky1976survey,isermann2006fdi}: her fiziksel kanal için kanal
        başına ridge regresyon --- girdi: komutun son 1 s'lik üçgen-ağırlıklı
        özeti + bağlam ($V$, $V^2$, uçuş fazı, $V\times$komut etkileşimi);
        hedef: tepkinin gelecek 0.5 s ortalaması. Ridge seçiminin gerekçesi
        performans değil \emph{teşhis edilebilirlik}tir: katsayılar fiziksel
        yorum taşır. Sabit kanatta 6 kanal (aileron$\to$roll hızı,
        elevator$\to$pitch hızı, rudder+roll$\to$yaw hızı,
        gaz$\to$hava hızı türevi, pitch+gaz$\to$tırmanma hızı, çapraz-iz hatası);
        çok-rotorda 4 kanal (duruş hedefi$\to$duruş hızı, motor PWM dağılımı,
        toplam PWM$\to$dikey ivme, konum hedefi$\to$hız tepkisi).
\end{enumerate}

\subsection{Karar katmanı}

Anlık skorlar doğrudan alarma çevrilmedi; literatürdeki değişim-noktası tespiti
çerçevesi kullanıldı \cite{page1954cusum,basseville1993detection}: kanal başına
çift yönlü Page CUSUM,
\[
S^{+}_{t} = \max\big(0,\; S^{+}_{t-1} + z_t - k\big), \qquad
S^{-}_{t} = \max\big(0,\; S^{-}_{t-1} - z_t - k\big),
\]
eşikler \emph{yalnız normal verideki doğal alarm yükünden} (blok-bootstrap ile)
kalibre edilip donduruldu; alarm sonrası 60 s yenileme (refractory) penceresi ve
olay-birleştirme uygulandı. K-of-N kalıcılık kuralları alternatif karar katmanı
olarak denendi.

\subsection{Sağlama (sanity) kapıları}
\label{sec:sanity}

Projenin ikinci yarısında, her aday skor kalibrasyona girmeden önce dört
eşikten-bağımsız sağlama testinden geçirildi. Bu kapılar, ilk yarıda yaşanan ve
\S\ref{sec:kn-artefakt}'te anlatılan artefaktın tekrarını önlemek için kondu:

\begin{enumerate}[itemsep=2pt]
  \item \textbf{Genlik testi:} skor ile ham girdi normu $\lVert x\rVert$
        arasındaki Spearman korelasyonu $\rho < 0.5$ olmalı (skor, genliğin
        kopyası olmamalı).
  \item \textbf{Eğitilmemiş-eş testi:} eğitilmiş modelin skor sıralaması, aynı
        mimarinin rastgele-ağırlıklı kopyasıyla $\rho < 0.7$--$0.8$ olmalı
        (eğitim gerçekten bir şey katmalı).
  \item \textbf{Başlangıç ayrışması (KS) testi:} etiketli olaylarda artık
        büyüklüğünün başlangıç-sonrası dağılımı, başlangıç-öncesinden
        Kolmogorov--Smirnov testiyle ayrışmalı ($p<0.01$). Bu test sinyalin
        \emph{varlığını} eşikten bağımsız ölçer; sinyal-yokluğu ile
        eşik-yanlışlığını birbirinden ayırır.
  \item \textbf{Komut-ablasyon testi:} komut girdileri çıkarılıp yalnız
        bağlamla eğitilen ``sakat'' model, tam modele yakın performans
        veriyorsa model komutu değil bir sızıntıyı öğreniyordur
        (varyans oranı $\geq 1.15$ şartı).
\end{enumerate}

\section{Deneysel Sonuçlar --- Veri Seti Bazında Ölçülmüş Skorlar}
\label{sec:sonuc}

\subsection{ALFA}
\label{sec:alfa-sonuc}

\begin{table}[h]
\centering\small
\caption{ALFA --- ana sonuçlar. FA = yanlış alarm.}
\label{tab:alfa-sonuclar}
\begin{tabular}{p{5.2cm}p{3.4cm}p{2.6cm}p{3.4cm}}
\toprule
Yöntem/kanal & Metrik & Değer & Sonuç \\
\midrule
Monolitik satır-bazlı Isolation Forest & ROC & 0.50 & Rastgele; başarısız \\
LSTM-AE (uçuş düzeyi) & uçuş-ROC & 0.731 & IF-füzyonun (0.833) altında \\
LSTM-AE / IF / LightGBM & AUPRC & \textbf{0.872} / 0.858 / 0.843 &
Literatürle uyumlu, iyi \\
İrtifa hatası + CUSUM (nedensel olmayan ilk ölçüm) & uçuş-ROC & 0.878 &
Yanıltıcı (aşağıya bkz.) \\
Aynı kanal, \emph{nedensel} yeniden ölçüm & uçuş-ROC & \textbf{0.611} &
``En güçlü sinyal'' sanılan kanal çöktü \\
Olay-başlangıcı yakalama, point-adjust benzeri şişirme düzeltmesi öncesi/sonrası &
recall & 0.594 $\to$ \textbf{0.194--0.224} & Literatürdeki şişirme etkisinin
kendi verimizdeki ölçümü \\
USAD (sürüklenme/sapma/donma enjeksiyonu) & recall & 0.45 / 0.53 & LSTM-AE'nin
altında; elendi \\
CUSUM, rudder sınıfı ($n{=}4$) & ``1.00 tespit'' & $n{=}1$'e dayalı &
İstatistiksel olarak anlamsız; raporlamadan çıkarıldı \\
\bottomrule
\end{tabular}
\end{table}

Ana sonuçlar Tablo~\ref{tab:alfa-sonuclar}'tedir. ALFA'nın öğrettiği iki şey:
(i) araştırma metriğinde literatür seviyesi
yakalanabiliyor (AUPRC 0.872); (ii) nedensellik ve olay-bazlı ölçüm
sıkılaştırıldığında görünür başarının büyük kısmı buharlaşıyor. Küçük sınıflar
($n{=}1$--$4$) üzerinde yüzde raporlamak reddedildi; bu sınıflar yalnız vaka
analiziyle ele alındı.

\subsection{UAV Attack}

Ana bulgular \S3.2'de verildi: ağ-katmanı DoS 6 logdan 4'ünde fiziksel
telemetride görünmez; 2 m/s sinsi GPS kayması artık uzayında normal
varyasyondan ayrılamaz; kalibrasyonsuz model transferi normal uçuşlarda
\textbf{1.00 yanlış alarm oranı} üretti (her normal uçuş alarm aldı). Monolitik
Isolation Forest bu sette ROC \textbf{0.21} verdi --- rastgeleden kötü; bunun
nedeni sonradan simülasyon/gerçek log karışımının sinyali sistematik ters
çevirmesi olarak teşhis edildi (\S\ref{sec:kn-domain}).

\subsection{UAV-SEAD}
\label{sec:sead-sonuc}

\begin{table}[h]
\centering\small
\caption{UAV-SEAD --- seçilmiş sonuçlar. Hedefler: uyarı $\geq 0.50$ recall @
$\leq 12$ FA/saat; kritik $\geq 0.30$ @ $\leq 2$.}
\label{tab:sead-sonuclar}
\begin{tabular}{p{5.6cm}p{3.2cm}p{3.2cm}p{2.6cm}}
\toprule
Yöntem & Recall & FA/uçuş-saati & Sonuç \\
\midrule
IF-füzyon, adil satır ölçümü & ROC 0.474 & --- & Sinyalsiz \\
EKF test-oranı kanalı & korelasyon \emph{ters} (0.354) & --- & Arızada ölçüm
reddi sinyali bastırıyor \\
LightGBM pencere & AUPRC 0.349 (IF 0.385) & --- & İyileşme yok \\
Dikey tutarlılık kanalı & $+0.021$ recall farkı & --- & Eşik altı ($<0.05$) \\
Motor simetrisi kanalı & $+0.024$ (2/5 tohum) & --- & Kararsız \\
Chronos sıfır-atış, motor kanalı & 0.205 $\to$ \textbf{0.390} & --- &
Gerçek yöntemsel iyileşme \\
İtki komutu kanalı (tek özellik) & 0.205 $\to$ \textbf{0.459} & \textbf{38.1} &
En iyi kanal; FA bütçenin 3 katı \\
Tüm-kanal füzyonları (3 ayrı tur) & 0.213--0.222 & 23.7--25.8 & Bütçe aşımı \\
İki-kanal ağırlıklı varyant & 0.217 $\to$ 0.291 & 23.70 $\to$ \textbf{44.70} &
Kazanım FA şişirilerek satın alındı \\
Normal havuz 899'a büyütüldükten sonra & 0.126 & 9.95 & Bütçeye girdi ama
recall çöktü \\
Derin modeller (LSTM-AE / Dense-AE / USAD), genlik düzeltmesi sonrası &
$<0.06$ & --- & Bkz.\ \S\ref{sec:kn-artefakt} \\
\bottomrule
\end{tabular}
\end{table}

Tablo~\ref{tab:sead-sonuclar}'daki örüntü tipiktir ve projenin ana duvarını
gösterir: \emph{her recall
kazancı, yanlış alarm bütçesini aşarak satın alınabiliyor; bütçeye girildiğinde
recall hedefin altına düşüyor.} Heterojen ve etkin-örneklemi küçük ($\sim$64
oturum) normal havuzda, yenilik-tespiti net bir ``normal'' referansı kuramıyor.

\subsection{RflyMAD}
\label{sec:rfly-sonuc}

Tablo~\ref{tab:rfly-sonuclar}, \S\ref{sec:rfly-veri}'te anlatılan etiket
granülerliği düzeltmesinin sayısal etkisini ve havuzlamanın maliyetini
gösterir: uçuş-düzeyi vekil etiketle parlak görünen sonuç, gerçek olay-aralığı
etiketiyle operasyonel bütçenin dışına düşer; iki veri setini havuzlamak ise
en iyi kanalı da bozar.

\begin{table}[h]
\centering\small
\caption{RflyMAD --- etiket düzeltmesi öncesi/sonrası.}
\label{tab:rfly-sonuclar}
\begin{tabular}{p{6cm}p{2.6cm}p{3cm}p{3cm}}
\toprule
Konfigürasyon & Recall & FA/uçuş-saati & Sonuç \\
\midrule
İtki komutu kanalı, uçuş-düzeyi \emph{vekil} etiket & 0.749 & --- &
\textbf{Geçersiz} (kolay soru) \\
Aynı kanal, gerçek \emph{aralık} etiketi (uyarı) & 0.526 & 22.28 & Recall iyi,
FA bütçenin $\sim$2 katı \\
Aynı kanal, kritik sözleşme & 0.442 & 9.23 & Kritik bütçe ($\leq 2$) aşıldı \\
İki set havuzlanınca (SEAD+RflyMAD) & \textbf{0.149} & 30.00 & Havuzlama ciddi
kötüleştiriyor \\
Havuzlanmış füzyon, bütçe içi varyant & 0.066 & 9.39 & Recall anlamsız düzeyde \\
\bottomrule
\end{tabular}
\end{table}

\subsection{ADS-B}
\label{sec:adsb-sonuc}

İlk tur, üç sinir ağını kural-bazlı fizik skorlayıcıyla aynı protokolde
yarıştırdı (Tablo~\ref{tab:adsb-auc}):

\begin{table}[h]
\centering\small
\caption{ADS-B --- sentetik bozulma senaryolarında AUC. Vekil pencere
etiketlemesinin yapısal tavanı $\sim$0.75'tir.}
\label{tab:adsb-auc}
\begin{tabular}{lcccc}
\toprule
Senaryo & Kural-bazlı & Dense-AE & LSTM-AE & LSTM öngörücü \\
\midrule
Havuzlanmış (tümü) & \textbf{0.600} & 0.572 & 0.568 & 0.552 \\
Yer hızı sapması & --- & 0.737 & 0.743 & 0.648 \\
Dikey hız donması & --- & 0.600 & 0.579 & 0.552 \\
Track donması & --- & 0.523 & 0.521 & 0.551 \\
Sinsi konum rampası & --- & 0.513 & 0.519 & 0.513 \\
İrtifa kesintisi & --- & 0.489 & 0.480 & 0.498 \\
\bottomrule
\end{tabular}
\end{table}

Üç sonuç öne çıkar: (i) \textbf{öğrenmesiz kural-bazlı skorlayıcı üç sinir ağını
da geçti}; (ii) sinir ağları yalnız büyük-genlikli bozulmayı (hız sapması)
yakaladı, sinsi senaryolarda rastgeleden ayrışamadı; (iii) üç sinir ağı da
genlik-baskınlığı bayrağı aldı (eğitilmiş--rastgele $\rho$ 0.84--0.94,
\S\ref{sec:kn-artefakt}).

Kural katmanı, düzeltilmiş olay-aralığı etiketleriyle yeniden ölçüldüğünde
güçlü ayrışma gösterdi: \textbf{AUROC 0.765, AUPRC 0.883}; yer hızı olaylarında
0.964 olay-yakalama (medyan gecikme 19.3 s), track olaylarında 0.952 (56.8 s).
Ancak operasyonel tablo yine kapanmadı:

\begin{itemize}[itemsep=2pt]
  \item Sinsi rampa senaryosunda olay-yakalama 0.801 \emph{görünse de}, olay
        süresi boyunca aktif kapsama yalnız \textbf{0.184} --- olaya bir kez
        dokunmak ile olay boyunca alarmda kalmak aynı şey değildir.
  \item Doğal (temiz) trafikte alarm yükü \textbf{4.81 olay/saat}; CUSUM karar
        katmanı eklendiğinde skorlanabilir uçuşların $\sim$\%99'u alarm aldı
        (alarm doygunluğu) --- operatör açısından kullanılamaz.
  \item Sonradan tasarlanan bağlamsal olasılıksal LSTM sağlıklı eğitildi (2.929
        uçuş, 1.27 milyon satır; Gauss NLL $0.795 \to 0.709$, $\sim$\%10.9
        düşüş; 9.546 parametre) ve genlik kapılarını geçti
        ($\rho = 0.650/0.654 < 0.8$); ancak karar katmanı dondurulamadan hat
        kapandı: doğal alarm doygunluğu aşılamadı ve dondurulmuş skor
        kaynağının bit-düzeyi yeniden üretimi sağlanamadığı için değerlendirme
        temkinli biçimde (fail-closed) durduruldu. Bu model için olay-yakalama
        sayısı raporlamak, ölçülmemiş bir sayı uydurmak olurdu; raporlanmadı.
\end{itemize}

\subsection{GNSS bütünlük alt-çalışması}
\label{sec:gnss-sonuc}

Üç yöntem (PX4-yerleşik yenilik/test-oranı kuralları; çok-kanallı CUSUM;
bağlamsal LSTM) iki rol (geliştirme, prova) ve iki sözleşme (kritik, uyarı)
üzerinde, kalibrasyonda dondurulan eşiklerle değerlendirildi. \textbf{On iki
konfigürasyonun tamamı başarısız oldu} (Tablo~\ref{tab:gnss-sonuclar}):

\begin{table}[h]
\centering\small
\caption{GNSS --- dondurulmuş eşiklerle sonuçlar. Bütçeler: kritik $\leq 2$,
uyarı $\leq 12$ alarm/saat; recall hedefi kritik/uyarı için zamanında yakalama.}
\label{tab:gnss-sonuclar}
\begin{tabular}{llrrl}
\toprule
Rol & Yöntem (sözleşme) & Recall & Alarm/saat & Başarısızlık nedeni \\
\midrule
Geliştirme & PX4-yerleşik (kritik) & \%58.8 & 19.58 & Alarm yükü $\sim$10 kat \\
Geliştirme & CUSUM (kritik) & \%0.0 & 0.00 & Bütçeyi sağlayan sonlu eşik yok \\
Geliştirme & LSTM (kritik) & \%47.1 & 20.32 & Alarm yükü \\
Geliştirme & PX4-yerleşik (uyarı) & \%58.8 & 19.58 & Alarm yükü \\
Geliştirme & CUSUM (uyarı) & \%52.9 & 0.00 & Recall hedefi altı \\
Geliştirme & LSTM (uyarı) & \%82.4 & \textbf{101.60} & Aşırı alarm yükü \\
Prova & PX4-yerleşik (kritik) & \%90.0 & 28.86 & Alarm yükü \\
Prova & CUSUM (kritik) & \%0.0 & 0.00 & Eşik sonsuz \\
Prova & LSTM (kritik) & \%90.0 & 0.00 & Geliştirmeye genellemiyor \\
Prova & PX4-yerleşik (uyarı) & \%90.0 & 28.86 & Alarm yükü \\
Prova & CUSUM (uyarı) & \%70.0 & 7.21 & \%90 hedefi altı \\
Prova & LSTM (uyarı) & \%100.0 & 44.86 & Alarm yükü \\
\bottomrule
\end{tabular}
\end{table}

Dikkat çeken örüntü LSTM satırlarındadır: prova kümesinde \%90 recall / sıfır
alarm gibi görünen model, \emph{aynı dondurulmuş kararla} geliştirme kümesinde
\%47.1 recall / 20.32 alarm/saat üretmiştir. Bu, kalibrasyon kümesine özgü
başarı ve dağılım kararsızlığıdır; ürün genellemesi değildir. Rüzgâr stres
testinde (886 SIL/HIL vakası) dondurulmuş dedektörler simülasyonda bütçe içinde
kalırken (8.27/6.88 alarm/saat) donanım-döngüde bütçeyi aştı
(\textbf{20.53/22.61} $> 12$): çevresel alan kaymasına aşırı hassasiyet.
Kritik CUSUM eşiğinin ``sonsuz'' kalması yazılım hatası değildir: kalibrasyon
kümesinde 2 alarm/saat bütçesini sağlayan sonlu eşik \emph{bulunamamış}, güvenli
davranış ``hiç alarm üretme'' olmuştur. Karar: \textbf{NO-GO --- mevcut veri ve
enstrümantasyonla ürün adayı elde edilemez}; kör test hiç açılmadı.

\subsection{Komut$\to$tepki artık modeli (son çalışma)}
\label{sec:residual-sonuc}

Projenin son çalışması, anomaliyi ham telemetride değil \emph{öğrenilmiş uçuş
dinamiği modelinin yeniliğinde} arayan klasik model-bazlı FDI çerçevesine
döndü. Kanal başına ridge sonuçları Tablo~\ref{tab:ridge}'dedir:

\begin{table}[h]
\centering\small
\caption{Kanal başına ridge regresyon --- çapraz doğrulama $R^2$.}
\label{tab:ridge}
\begin{tabular}{p{5.6cm}rrp{4.6cm}}
\toprule
Kanal & CV $R^2$ & Train $R^2$ & Durum \\
\midrule
Sabit kanat: 5 komut$\to$tepki kanalı & --- & --- & Eğitilemedi (bölme için
$\geq 2$ oturum gerekir; 1 var) \\
Sabit kanat: çapraz-iz hatası & \multicolumn{2}{c}{öğrenmesiz kanal} &
Kullanıldı \\
Çok-rotor: duruş hedefi$\to$duruş hızı & 0.011 & 0.364 & Zayıf; aşırı uyum \\
Çok-rotor: motor PWM dağılımı & \textbf{0.456} & 0.711 & En güçlü öğrenilmiş
kanal \\
Çok-rotor: toplam PWM$\to$dikey ivme & 0.000 & 0.009 & Sinyalsiz \\
Çok-rotor: konum hedefi$\to$hız tepkisi & --- & --- & Uygun satır yok \\
\bottomrule
\end{tabular}
\end{table}

Sağlama kapıları (\S\ref{sec:sanity}) bu kez \emph{tasarımın parçasıydı} ve
çalıştı: komut-ablasyonu iki kanalın komutu gerçekten kullandığını doğruladı
(varyans oranları 1.199 ve 2.497), üçüncü kanalı elemeye yaradı (1.008);
genlik testi üç aktif kanalda geçti ($\rho = 0.472, 0.140, 0.018 < 0.5$).
En önemlisi, \textbf{başlangıç-ayrışması (KS) testi üç başlık sınıfının
üçünde de geçti}:

\begin{center}
\small
\setlength{\tabcolsep}{4pt}
\begin{tabular}{lcccc}
\toprule
Sınıf / kanal & KS & $p$ & Öncesi medyan $|z|$ & Sonrası medyan $|z|$ \\
\midrule
Sabit kanat motor / çapraz-iz & 0.165 & $5.5\times10^{-18}$ & 1.13 & 1.81 \\
Çok-rotor motor / duruş kanalı & 0.177 & $\approx 0$ & 0.98 & 1.50 \\
Çok-rotor motor / motor PWM dağılımı & \textbf{0.570} & $\approx 0$ & 1.10 &
\textbf{5.79} \\
Çok-rotor sensör / duruş kanalı & 0.274 & $\approx 0$ & 0.88 & 1.59 \\
Çok-rotor sensör / motor PWM dağılımı & 0.034 & $2.3\times10^{-91}$ & 0.95 &
1.02 \\
\bottomrule
\end{tabular}
\end{center}

Yani \emph{sinyal eşikten bağımsız olarak vardır}. Buna rağmen çalışma NO-GO
ile kapandı, çünkü eşik kalibrasyonu aşamasında şu ölçüldü: hedef yanlış alarm
oranında tek bir alarmı dahi güvenilir çözebilmek için gereken bağımsız normal
uçuş maruziyeti mevcut değil:

\begin{center}
\small
\begin{tabular}{lrrr}
\toprule
Veri & Mevcut normal maruziyet & Gerekli asgari & Açık \\
\midrule
Sabit kanat (ALFA) & 0.169 saat & 2.0 saat & \textbf{11.8$\times$} \\
Çok-rotor (RflyMAD) & 0.786 saat & 4.0 saat & \textbf{5.1$\times$} \\
\bottomrule
\end{tabular}
\end{center}

Bootstrap mevcut blokları yeniden örnekleyebilir ama \emph{yeni bağımsız
uçuş-saati yaratamaz}. Veri tavanı analizi de kapıyı kapattı: ALFA'nın resmî
külliyatı 47 uçuşta sabittir (11 normal; en iyimser yeniden dağıtım $+\%22$
sağlar, 11.8$\times$ açığı kapatmaz); RflyMAD'in resmî arızasız tavanı 84'tür
(projede 51; kalan en iyimser 33 aday eklense bile 1.42 saate ulaşılır, 4.0
saat hedefinin 2.58 saat altında --- açığı kapatmak $\sim$135 ek normal uçuş
gerektirir). Eşikler yazılmadı, kör test açılmadı, yanlış alarm hedefi sonuca
bakılarak gevşetilmedi.

\section{Kök Neden Analizi: Proje Neden Başarısız Oldu?}
\label{sec:kokneden}

Yukarıdaki sonuçlar yedi kök nedene indirgenebilir. Her biri için kanıt,
önceki bölümlerin ölçülmüş sayılarıdır.

\subsection{KN-1: Değerlendirme çıtası uyumsuzluğu --- literatürün ölçmediği
bir büyüklüğe hedef kondu}

Literatürdeki ``\%90+'' sonuçların büyük çoğunluğu AUROC/AUPRC/F1 ailesindendir
ve dondurulmuş eşikte saatlik yanlış alarm yükü \emph{genellikle hiç
raporlanmaz} \cite{wu2021flawed,kim2022rigorous}. Proje bu metriklerde
literatürü yakalamıştır (ALFA AUPRC 0.872; ADS-B kural katmanı AUROC 0.765 /
AUPRC 0.883). Takılınan nokta, literatürün ölçmediği operasyonel çift metrik
oldu. Bu bir teselli cümlesi değil, planlama hatasının tespitidir:
\emph{operasyonel bir çıta, ancak o çıtayı istatistiksel olarak taşıyabilecek
veri kampanyasıyla birlikte anlamlıdır.} Araştırma-kalitesi küçük veriyle
operasyonel-dağıtım seviyesi bir iddia kanıtlanmaya çalışıldı.

\subsection{KN-2: Normal uçuş maruziyeti yapısal olarak yetersiz}

Saatte $\leq 2$ yanlış alarm gibi bir bütçeyi kalibre etmek, tanımı gereği çok
sayıda bağımsız normal uçuş-saati gerektirir. Mevcut havuzlar: ALFA 11 normal
uçuş ($0.169$ saat kullanılabilir maruziyet), UAV Attack 6 log, UAV-SEAD
$\sim$64 bağımsız oturum, RflyMAD 51 arızasız uçuş ($0.786$ saat). Son
çalışmanın kalibrasyon analizi açığı kesinleştirdi: gereken maruziyetin
$1/11.8$'i ve $1/5.1$'i mevcuttu ve resmî külliyat tavanları (47 / 84 uçuş) bu
açığı \emph{matematiksel olarak} kapatamıyordu. Etiketli akademik setler
``hangi uçuş arızalı'' sorusu için tasarlanmıştır; ``normal uçuşta saatte kaç
kez boşa öteriz'' sorusu için değil.

\subsection{KN-3: Heterojen normal havuzunda yenilik-tespitinin çökmesi}

``Yalnız normali öğren, sapanı işaretle'' yaklaşımı, normalin kendisi
dağınıksa bozulur; anomali tespiti literatürü bu olguyu (swamping/masking,
normallik tanımının belirsizliği) bilinen bir zorluk olarak kaydeder
\cite{ruff2021unifying}. Dört etiketli sette de normal havuz hem küçük hem
heterojendi. Ölçülen belirtiler: her recall kazanımının FA şişmesiyle gelmesi
(UAV-SEAD: recall $+0.074$ için FA $1.89\times$; K-of-N varyantında
$3.83\times$), normal havuz 899 uçuşa büyütüldüğünde recall'un 0.126'ya
çökmesi (veri artışı, heterojenliği de artırdı), iki seti havuzlamanın en iyi
kanalı 0.526'dan 0.149'a düşürmesi.

\subsection{KN-4: Genlik-baskınlığı artefaktı --- yeniden-kurma tabanlı
modellerin gerçek teknik hatası}
\label{sec:kn-artefakt}

Projenin en önemli \emph{teknik} bulgusu şudur: üç farklı derin mimari
(LSTM-AE, Dense-AE, USAD) UAV-SEAD üzerinde neredeyse aynı skorları üretti ve
teşhis testleri nedenini gösterdi --- eğitilmiş modellerin skor sıralaması,
aynı mimarinin \emph{rastgele-ağırlıklı} kopyasıyla $\rho = 0.964$, ham girdi
enerjisiyle ($\lVert x\rVert^2$) $\rho = 0.965$ korelasyondaydı. Yani modeller
örüntü değil, \textbf{genlik} öğrenmişti; kırpılmamış sağlam-ölçekleme, bir
avuç aşırı-genlikli pencerenin skoru domine etmesine izin vermişti. Göreli-hata
düzeltmesi uygulandığında korelasyonlar kırıldı ($\rho = 0.13$--$0.23$) ama
altından operasyonel sinyal çıkmadı: recall $<0.06$. ADS-B'deki üç sinir ağı da
aynı bayrağı aldı ($\rho = 0.84$--$0.94$).

Bu artefaktın literatürdeki karşılığı, sığ/derin anomali yöntemlerinin
``kestirme öğrenme'' (shortcut learning) davranışı ve kusurlu kıyaslarda
önemsiz çözülebilirlik bulgusudur \cite{wu2021flawed,ruff2021unifying}.
Projede bu hata \emph{yakalandı, raporlandı ve sonraki tüm modellere zorunlu
sağlama kapısı yapıldı} (\S\ref{sec:sanity}); nitekim sonradan tasarlanan
bağlamsal LSTM ($\rho = 0.650$) ve komut$\to$tepki artıkları ($\rho \leq
0.472$) bu kapıları geçti. Ancak ilk yarıdaki derin-model sonuçlarının
tamamını geçersizleştirmesi, projeye ciddi zaman kaybettirdi.

\subsection{KN-5: Etiket granülerliği uyumsuzluğu --- kolay soruya doğru cevap}

Uçuş-düzeyi etiketle ölçmek (``bu uçuşta arıza var mı'') ile olay-aralığı
etiketiyle ölçmek (``arıza hangi saniyede başladı, ne kadarını kapsadın'')
arasındaki fark projede iki kez sayısallaştı: RflyMAD'de 0.749 $\to$ 0.526;
ALFA'da point-adjust benzeri şişirme düzeltmesiyle 0.594 $\to$ 0.194--0.224.
Literatür bu şişirmeyi sistematik olarak belgelemiştir \cite{kim2022rigorous}.
ADS-B tarafında vekil pencere etiketlemesinin $\sim$0.75'lik yapısal AUC tavanı
da aynı ailedendir: \emph{etiketin sorduğu soru ile operasyonel soru aynı
değilse, metrik ne kadar parlaksa yorum o kadar yanıltıcıdır.}

\subsection{KN-6: Fiziksel gözlenebilirlik boşlukları --- sinyal sensörde yoksa
modelde de yoktur}

Dört ayrı ölçüm, bazı anomali türlerinin mevcut kolonlarda \emph{ilkesel
olarak} görünmez olduğunu gösterdi: (i) ağ-katmanı DoS, uçuş fiziği kolonlarına
yansımıyor (6 logdan 4'ünde tespit edilemez); (ii) RflyMAD'de ESC telemetrisi
(RPM/akım/voltaj/sıcaklık) şemada var ama tüm uçuşlarda sabit sıfır --- motor
sağlığı pilotu bu yüzden hiç açılamadı; (iii) ADS-B'de barometrik--geometrik
irtifa farkı kanalının fit MAD'ı tam sıfır; (iv) 2 m/s sinsi GPS kayması artık
uzayında normal varyasyon içinde kalıyor ve sinsi rampada olay-içi aktif
kapsama 0.184'te tıkanıyor. Bu vakalarda başarısızlık modelin değil,
\emph{enstrümantasyonun} sınırıdır; hiçbir yöntem ailesi bu boşluğu
kapatamazdı.

\subsection{KN-7: Alan (domain) kayması ve genellemesizlik}
\label{sec:kn-domain}

Simülasyon/gerçek karışımı UAV Attack'te sinyali sistematik ters çevirdi
(ROC 0.21; EKF test-oranı kanalında ters korelasyon --- arıza sırasında ölçüm
reddi mekanizması yenilik sinyalini bastırıyor). GNSS çalışmasında aynı
dondurulmuş LSTM kararı prova kümesinde \%90 / sıfır alarm, geliştirme
kümesinde \%47.1 / 20.32 alarm-saat üretti; rüzgâr stresinde simülasyon bütçe
içi, donanım-döngü bütçe dışı (20.53/22.61 $>$ 12) çıktı. İki veri setini
havuzlamak en iyi kanalı 0.526'dan 0.149'a düşürdü. Küçük veri + heterojen
koşullar altında öğrenilen karar sınırları roller/alanlar arasında taşınmadı.

\subsection{Kök nedenlerin etkileşimi}

Bu yedi neden bağımsız değildir; birbirini besledi. Küçük normal havuz (KN-2)
heterojenliği görünmez kıldı (KN-3); heterojenlik derin modelleri genlik
kestirmesine itti (KN-4); parlak görünen uçuş-düzeyi metrikler (KN-5) bu
artefaktın geç fark edilmesine yol açtı; operasyonel çıta (KN-1) ise tüm bu
zincirin en zayıf halkasını --- kalibrasyon maruziyetini --- sürekli test etti
ve her seferinde kırdı. Son iki çalışmanın katkısı, bu zinciri ayrıştırıp
``sinyal var mı'' (evet; KS testleri) ile ``eşik kurulabilir mi'' (hayır;
maruziyet açığı) sorularını ayrı ayrı, eşikten bağımsız biçimde
cevaplamasıdır.

\section{Neler Çalıştı --- Dürüst Pozitifler}

Başarısızlık raporunun simetrik yükümlülüğü, neyin çalıştığını da kaydetmektir:

\begin{itemize}[itemsep=2pt]
  \item \textbf{Fizik-tutarlılık artıkları + kural katmanı}, ADS-B'de üç sinir
        ağını da geçti (AUC 0.600 > 0.552--0.572) ve düzeltilmiş etiketlerle
        AUROC 0.765 / AUPRC 0.883'e ulaştı; bariz senaryolarda olay-yakalama
        0.95+ idi. Basit ve teşhis edilebilir yöntemin karmaşık olanı yendiği
        bu bulgu, literatürdeki ``tek kazanan yok'' sonucuyla tutarlıdır
        \cite{schmidl2022eval}.
  \item \textbf{Sıfır-atış temel model} (Chronos), alan-içi eğitim olmadan
        motor kanalında 0.205 $\to$ 0.390 iyileşme sağladı --- projedeki nadir
        gerçek yöntemsel kazanımlardan biri \cite{ansari2024chronos}.
  \item \textbf{Sağlama kapıları mimarisi} (genlik, eğitilmemiş-eş, ablasyon,
        KS) ölümcül bir artefaktı yakaladı ve tekrarını önledi; sonraki
        modeller bu kapılardan geçerek eğitildi.
  \item \textbf{Model-bazlı FDI dönüşü} sinyalin varlığını eşikten bağımsız
        kanıtladı (motor PWM dağılım kanalında KS 0.570; başlangıç sonrası
        medyan $|z|$ 1.10 $\to$ 5.79) ve en güçlü öğrenilmiş kanalı üretti
        (CV $R^2$ 0.456).
  \item \textbf{Metodolojik disiplin}: kör test kümeleri hiç açılmadı; eşikler
        sonuca bakılarak gevşetilmedi; şüpheli etiketli uçuşlar sonuç
        görülmeden dışlandı; her NO-GO sayısal gerekçeyle belgelendi.
        Sonuç tablosundaki ``başarısız'' satırlarının bir kısmı, hile yapan bir
        değerlendirmenin parlak göstereceği konfigürasyonlardır.
\end{itemize}

\section{Dersler ve Öneriler}

\begin{enumerate}[itemsep=3pt]
  \item \textbf{Operasyonel çıta, veri kampanyasıyla birlikte tasarlanmalı.}
        ``$\leq \lambda$ yanlış alarm/saat'' iddiası için gereken bağımsız
        normal uçuş-saati \emph{deney başlamadan} hesaplanabilir ve
        hesaplanmalıdır; mevcut açık 5--12 kat idi. Bu ölçekte açık, yöntem
        seçimiyle kapanmaz.
  \item \textbf{Olay-aralığı etiketi olmayan veri setiyle operasyonel iddia
        kurulmamalı.} Uçuş-düzeyi etiket yalnız araştırma metriği taşır;
        point-adjust ve vekil etiket şişirmeleri kendi verimizde 2--3 kat
        fark yarattı.
  \item \textbf{Sağlama kapıları en başa konmalı.} Genlik ve eğitilmemiş-eş
        testleri ilk turda mevcut olsaydı, derin-model turlarının tamamına
        yakını daha ilk hafta geçersiz sayılacaktı.
  \item \textbf{Yeniden-kurma tabanlı otokodlayıcılar, manevralı İHA
        telemetrisi için uygun varsayılan değil.} Manevra genliği ile anomali
        genliği aynı eksende yarışıyor; komut$\to$tepki artığı (model-bazlı
        FDI) manevrayı yapısal olarak skordan düşürüyor. Devam edilecekse
        buradan devam edilmeli.
  \item \textbf{Gözlenebilirlik denetimi model işinden önce gelir.} Kolon
        şemasının varlığı, sinyalin varlığı değildir (sıfır-varyanslı ESC
        alanları; telemetriye yansımayan ağ saldırıları). Her hedef anomali
        türü için ``hangi kolonda, hangi fizikle görünür?'' sorusu yazılı
        cevaplanmalı.
  \item \textbf{İnterpolasyon Silver katmanına sokulmamalı}; her kanal doğal
        hızında tutulup geçmişe-bakan hizalama ve bayatlık kolonlarıyla
        işlenmeli.
  \item \textbf{Yeniden açma koşulları} (her ikisi de veri kararıdır, yöntem
        kararı değil): (a) kontrollü yeni normal-uçuş kampanyası (çok-rotor
        için $\sim$135+ ek normal uçuş ölçeğinde) ve fiziksel olarak
        belgelenmiş arıza enjeksiyonu; veya (b) yanlış alarm bütçesinin ayrı
        bir ön-kayıtla, yeni bir deney sözleşmesi olarak yeniden müzakeresi.
\end{enumerate}

\section{Sonuç}

Proje, hedeflediği operasyonel anomali dedektörünü üretememiştir; bu rapor
bunun nedenlerini sayı sayı belgelemiştir. Bununla birlikte nihai tablo
``hiçbir şey öğrenilemedi'' değildir: eşikten bağımsız testler en az üç
sınıfta gerçek sinyal kanıtlamış; genlik-baskınlığı artefaktı gibi literatürde
de yaygın bir değerlendirme tuzağı yakalanıp kapatılmış; ve başarısızlığın ana
kaynağının yöntem ailesi değil, \emph{operasyonel çıta ile akademik veri
setlerinin taşıma kapasitesi arasındaki yapısal uyumsuzluk} olduğu
gösterilmiştir. Projenin son hâli iki cümleyle özetlenir: \textbf{sinyal var;
mevcut veriyle operasyonel karar sınırı kurulamaz.} Bu ayrımın kendisi ---
ve onu kuran eşikten-bağımsız sağlama/kalibrasyon metodolojisi --- projenin
aktarılabilir ana çıktısıdır.

\begin{thebibliography}{99}\small

\bibitem{keipour2021alfa}
A.~Keipour, M.~Mousaei, S.~Scherer,
``ALFA: A Dataset for UAV Fault and Anomaly Detection,''
\emph{The International Journal of Robotics Research}, 40(2--3), 2021.
arXiv:1907.06268.

\bibitem{whelan2020uav}
J.~Whelan, T.~Sangarapillai, O.~Minawi, A.~Almehmadi, K.~El-Khatib,
``Novelty-based Intrusion Detection of Sensor Attacks on Unmanned Aerial
Vehicles,'' \emph{Proc.\ ACM Q2SWinet}, 2020 (UAV Attack Dataset, IEEE
DataPort).

\bibitem{rflymad}
Rfly-OpenHA proje ekibi (Beihang Üniversitesi),
``RflyMAD: A Dataset for Multicopter Fault Detection and Health Assessment,''
resmî veri dokümantasyonu:
\url{https://rfly-openha.github.io/documents/4_resources/dataset.html}.

\bibitem{adsblol}
adsb.lol topluluk ADS-B ağı, günlük tam-ağ arşivleri.
\url{https://www.adsb.lol}.

\bibitem{liu2008iforest}
F.~T.~Liu, K.~M.~Ting, Z.-H.~Zhou,
``Isolation Forest,'' \emph{Proc.\ IEEE ICDM}, 2008.

\bibitem{ke2017lightgbm}
G.~Ke \etal{},
``LightGBM: A Highly Efficient Gradient Boosting Decision Tree,''
\emph{NeurIPS}, 2017.

\bibitem{malhotra2016lstm}
P.~Malhotra, A.~Ramakrishnan, G.~Anand, L.~Vig, P.~Agarwal, G.~Shroff,
``LSTM-based Encoder-Decoder for Multi-sensor Anomaly Detection,''
\emph{ICML Anomaly Detection Workshop}, 2016. arXiv:1607.00148.

\bibitem{audibert2020usad}
J.~Audibert, P.~Michiardi, F.~Guyard, S.~Marti, M.~A.~Zuluaga,
``USAD: UnSupervised Anomaly Detection on Multivariate Time Series,''
\emph{Proc.\ ACM KDD}, 2020.

\bibitem{ansari2024chronos}
A.~F.~Ansari \etal{},
``Chronos: Learning the Language of Time Series,''
\emph{Transactions on Machine Learning Research}, 2024. arXiv:2403.07815.

\bibitem{page1954cusum}
E.~S.~Page, ``Continuous Inspection Schemes,''
\emph{Biometrika}, 41(1--2), 1954.

\bibitem{basseville1993detection}
M.~Basseville, I.~V.~Nikiforov,
\emph{Detection of Abrupt Changes: Theory and Application},
Prentice Hall, 1993.

\bibitem{vovk2005conformal}
V.~Vovk, A.~Gammerman, G.~Shafer,
\emph{Algorithmic Learning in a Random World}, Springer, 2005.

\bibitem{willsky1976survey}
A.~S.~Willsky,
``A Survey of Design Methods for Failure Detection in Dynamic Systems,''
\emph{Automatica}, 12(6), 1976.

\bibitem{isermann2006fdi}
R.~Isermann, \emph{Fault-Diagnosis Systems: An Introduction from Fault
Detection to Fault Tolerance}, Springer, 2006.

\bibitem{kim2022rigorous}
S.~Kim, K.~Choi, H.-S.~Choi, B.~Lee, S.~Yoon,
``Towards a Rigorous Evaluation of Time-Series Anomaly Detection,''
\emph{Proc.\ AAAI}, 2022.

\bibitem{wu2021flawed}
R.~Wu, E.~Keogh,
``Current Time Series Anomaly Detection Benchmarks are Flawed and are
Misleading the Community,''
\emph{IEEE Transactions on Knowledge and Data Engineering}, 2021.

\bibitem{ruff2021unifying}
L.~Ruff \etal{},
``A Unifying Review of Deep and Shallow Anomaly Detection,''
\emph{Proceedings of the IEEE}, 109(5), 2021.

\bibitem{schmidl2022eval}
S.~Schmidl, P.~Wenig, T.~Papenbrock,
``Anomaly Detection in Time Series: A Comprehensive Evaluation,''
\emph{Proceedings of the VLDB Endowment}, 15(9), 2022.

\end{thebibliography}

\end{document}
